% Insertable results section. Requires amsmath, graphicx, booktabs and hyperref.
% Set \ADRResultsPath to this artifact directory before \input if needed.
\providecommand{\ADRResultsPath}{docs/research/solver_studies/adr_scaling_2026_09_17}
\providecommand{\ADRResultsFigurePath}{run_outputs/solver_studies/adr_scaling_2026_09_17/figures}
\section{GPU solver performance for stationary ADR}
\label{sec:adr-gpu-results}

We compare block additive Schwarz (ASM) and block Jacobi (BJ), augmented
by polynomial preconditioning, against block algebraic multigrid and
$p$MG--AMG: face-block modal polynomial multigrid with an AMGX coarse
hierarchy. Polynomial coarsening reduces the face degree on the same mesh;
AMGX algebraically coarsens the scalar degree-zero system. Configuration
labels PP($d$) report the polynomial degree, separately from the finite-element degree.
The smooth problems establish scaling; oscillatory transport and nonconvex
geometry expose robustness limits. $p$MG--AMG results now cover every matrix
and right-hand side in the study. Fresh and reused costs remain distinct,
and only physically converged solves enter timing comparisons.

\input{\ADRResultsPath/pardiso_lu/protocol.tex}

\subsection{Smooth three-regime scaling}
\label{sec:adr-smooth-scaling}
\paragraph{Problem and common protocol.}
We solve a manufactured stationary problem on $\Omega=[-1,1]^2$,
\begin{equation}
 \begin{aligned}
 -\Delta u+\boldsymbol\beta\cdot\nabla u+u&=f,
 &u|_{\partial\Omega}&=u_\star &&\text{(high diffusion)},
 \end{aligned}
\end{equation}
\begin{equation}
 \begin{aligned}
 -10^{-3}\Delta u+\boldsymbol\beta\cdot\nabla u+u&=f,
 &u|_{\partial\Omega}&=u_\star &&\text{(transport dominated)},
 \end{aligned}
\end{equation}
\begin{equation}
 \begin{aligned}
 -\nabla\cdot\!\left(\operatorname{diag}(1,10^{-2})\nabla u\right)
 +\boldsymbol\beta\cdot\nabla u+u&=f,
 &u|_{\partial\Omega}&=u_\star &&\text{(anisotropic)},
 \end{aligned}
\end{equation}
where $u_\star=\sin(\pi x)\sin(\pi y)$ and
$\boldsymbol\beta=(1,1/2)^T$ in all three cases.
On a grid of $n\times n$ split squares,
$h=2/n$, $N_T=2n^2$, and the condensed trace dimension is
$N_\lambda=(3n^2-2n)(p+1)$.
We vary $n\in\{32,64,128,223\}$ at $p=6$, and
$p\in\{1,2,3,4,6\}$ at $n=64$.
The largest system has $99\,458$ triangles and $1\,041\,187$ trace unknowns.

All methods receive the same FP64 Bernstein trace matrix and right-hand
side, with no outer scaling and a zero initial guess.
The HDG stabilization is $\tau=1+\max(\boldsymbol\beta\!\cdot\boldsymbol n,0)$.
\textbf{Acceptance is based on the independently recomputed physical residual}
$\|b-A\lambda\|_2/\|b\|_2\le10^{-10}$;
the common internal tolerance is $10^{-11}$, GMRES restart is 75, and the
iteration limit is 1,000. GPU measurements use one RTX PRO 5000 Blackwell GPU.
One setup is discarded as warmup, then three fresh setups each perform two
zero-guess solves.  Reported hot times average the second solve of each
setup; error bars span those three measurements.
Hot solve time covers only the linear solve; preconditioner setup and
host residual/reconstruction checks are reported separately or untimed.

\textbf{Configurations are frozen along each curve.}
ASM+PP(24) is used for high diffusion and transport, and ASM+PP(48)
for anisotropy, all with CGS2 GMRES.  AMGX uses natural $(p+1)\times(p+1)$ BSR blocks and a classical
hierarchy built from dense block adjacency, with block-Jacobi smoothing
for high/anisotropic diffusion and
block-DILU smoothing for transport.  Its FGMRES and BiCGSTAB runs share the
same preconditioner parameters.
Direct block-DILU is additionally tested
for transport.  $p$MG--AMG is tested in all three regimes: its hierarchy
preconditions the symmetric part in modal coordinates, with an AMGX scalar
coarse solve, while outer GMRES applies the original ADR operator.
The $p$MG--AMG standard policy uses direct degree-zero transfer, order-two
Chebyshev smoothing and one pre/post sweep. The robust policy uses successive
degree halving, order-four smoothing and two pre/post sweeps. Both policies
are shown wherever measured; the other configurations remain fixed along
each curve.

\subsubsection{Mesh and polynomial-order scaling}
\begin{figure}[!htbp]
 \centering
 \includegraphics[width=\linewidth]{\ADRResultsFigurePath/h_scaling.pdf}
 \caption{Mesh refinement at $p=6$: mean hot solve time (top) and Krylov
 iterations (bottom). LU appears in timing panels, with MKL thread counts
 beside its points. AMG and DILU curves use AMGX; $p$MG--AMG denotes the face-block modal hierarchy. Iterations are not equal units of work: BiCGSTAB usually applies
 its preconditioner twice per iteration.}
 \label{fig:adr-h-scaling}
\end{figure}

For high diffusion, polynomially preconditioned ASM grows from 33 to 150 iterations over the mesh ladder, whereas $p$MG--AMG grows from 58 to 71. At the largest size, $p$MG--AMG takes 0.278 s and the faster block-AMG variant 0.250 s, versus 3.469 s for polynomially preconditioned ASM. The small-mesh ranking therefore does not predict the large-mesh winner.
The direct-DILU comparison also separates the effect of the outer Krylov
method from that of a multilevel preconditioner: on the largest transport system, FGMRES/DILU takes 2.104 s (964 iterations), while BiCGSTAB/DILU takes 0.638 s (264 iterations).

\begin{figure}[!htbp]
 \centering
 \includegraphics[width=\linewidth]{\ADRResultsFigurePath/p_scaling.pdf}
 \caption{Polynomial-order scaling on $8\,192$ triangles. LU point labels
 give the MKL thread count. The block size
 grows from 2 to 7; trace dimension grows from $24\,320$ to $85\,120$.
 GPU parameters are fixed within each regime, but the algebraic hierarchies
 change with the matrix.  Thus these empirical curves need not be
 monotonic in $p$.}
 \label{fig:adr-p-scaling}
\end{figure}

For high diffusion at fixed mesh, $p$MG--AMG rises only from 39.8 to 55.2 ms as $p$ increases from 1 to 6 (50 to 65 iterations). AMGX and ASM exhibit nonmonotonic degree dependence under the frozen configurations; no universal power law is inferred.
The accuracy gain is shown in Figure~\ref{fig:adr-cost-accuracy};
$p$-refinement should therefore be judged at comparable error, not solely
by the time of one linear solve.

\subsubsection{Cost, robustness, and practical choice}
\begin{table}[!htbp]
 \centering
 \caption{Selected passing configurations on $99\,458$ triangles at $p=6$.
 Hot times are means; setup times are medians, in milliseconds. LU uses
 its reused-selected thread count for both columns.
 The block-AMG row in each regime is the faster of its two tested outer
 Krylov methods.}
 \label{tab:adr-largest}
 {\small\setlength{\tabcolsep}{3.3pt}
 \input{\ADRResultsPath/largest_table.tex}}
\end{table}

\begin{figure}[!htbp]
 \centering
 \includegraphics[width=\linewidth]{\ADRResultsFigurePath/cost_accuracy.pdf}
 \caption{Left: setup plus one hot solve for the selected large-system
 candidates (an arithmetic cost model, not a measured full PDE pipeline).
 F/B denote FGMRES/BiCGSTAB; High/Low/Aniso. identify diffusion regimes.
 For $R$ reused right-hand sides, the per-solve model is
 $T_{\rm hot}+T_{\rm setup}/R$.
 LU uses one thread count for both terms, shown in its label.
 Right: reconstructed primal $L^2$ error at fixed $8\,192$ triangles,
 using independent CPU direct-reference solutions.}
 \label{fig:adr-cost-accuracy}
\end{figure}

At the largest size, the fastest tested hot-solve configurations are AMG / BiCGSTAB for high diffusion (0.250 s), AMG / FGMRES for transport dominated (0.569 s), AMG / BiCGSTAB for anisotropic diffusion (0.622 s). Setup can change the preferred method for a single right-hand side; Table~\ref{tab:adr-largest} and Figure~\ref{fig:adr-cost-accuracy} keep it separate.
Prior application profiling on these same $p=6$ matrices attributes
roughly 92--96\% of polynomially preconditioned ASM GPU operation time to preconditioning:
one complete degree-24/48 application costs about 21/42 ms.
The $p$MG--AMG application costs about 2.6 ms; AMGX block-AMG costs depend
on the smoother (about 4--15 ms per cycle).
These separately instrumented numbers explain cost; they are not used to
rank the uninstrumented solves. Earlier profiling sampled solver-memory
increases of roughly 1.75 GiB for $p$MG--AMG and 2.8--3.3 GiB for the compared
ASM/FGMRES block-AMG paths.

\textbf{A stopping-trigger effect matters.}
With unchanged matrix and polynomial preconditioner, transport polynomially preconditioned ASM
at internal tolerance $5\times10^{-11}$ takes two 31-step restart cycles
(1.414 s); tightening to $10^{-11}$ retains one 32-step cycle (about 0.730 s).
The early projected-residual test had triggered a restart before the true
residual passed.  The scaling study uses the latter internal tolerance for
\emph{every} method; it does not mix the two timing protocols.

\textbf{Validation and scope.}
The original campaign attempts 96 solver/system combinations on 24 distinct systems; 96 pass and 0 fail. The 768 warmup/measured solves belonging to passing configurations have a worst physical relative residual of $1.69\times10^{-11}$. 18 systems also pass independent CPU direct-solution checks. Saved trace solutions agree across solvers within $8.25\times10^{-12}$ relative to the polynomially preconditioned ASM solution.
Failures are retained explicitly; none contribute to time rankings.
These reference curves use a smooth manufactured stationary solution.
The focused comparisons below add coefficient variation and isolate the
AMGX--polynomially preconditioned ASM crossover before the final oscillatory, nonconvex problem.
The combined report audits the $p$MG--AMG extension separately.
Configurations are selected from the preceding sweep; conclusions concern
the tested hardware and parameter range.  Exact AMGX internal multilevel
operator counts are unavailable, so iteration counts do not represent equal
work across methods.
